%% kapa: exam=1 round=24 year=2013 period=1 form=A subject=economics subject_ko=경제원론 seq=6 number=46
%% source: Q-Net 감정평가사 기출문제 (cst00309) artlSeq=5073802 file=extracted/1008/5073802/2013년 제24회 감정평가사 1차시험 기출문제/1교시(A형).hwp
%% conversion: kordoc 4.24.51 -> tools/gen1.py (auto); review=proofread-sonnet; flags=none
\begin{question}{46}
세 사람 A, B, C로 이루어진 어떤 경제에서 공공재에 대한 세 사람의 수요함수($Q_{ A }$, $Q_{ B }$, $Q_{ C }$ )는 각각 $Q_{ A } =10-P_{ A }$ , $Q_{ B } =10- \frac{ 1 } { 3 } P_{ B }$ , $Q_{ C } =5- \frac{ 1 } { 2 } P_{ C }$ 이고, 공공재의 한계비용은 20으로 일정할 때, 사회적 후생을 극대화시키는 공공재 생산량은? (단, $P_{ A }$ , $P_{ B }$ , $P_{ C }$ 는 A, B, C가 공공재에 지불하는 가격이다.)
\begin{choices}
  \item 5
  \item 10
  \item 15
  \item 20
  \item 25
\end{choices}
\end{question}
\answer{1}
